\documentclass[11pt]{article}

\usepackage[T1]{fontenc}
\usepackage[utf8]{inputenc}
\usepackage{amsmath,amssymb,amsthm,mathtools}
\usepackage[a4paper,margin=30mm]{geometry}
\usepackage{microtype}
\usepackage{xurl}
\usepackage[hidelinks]{hyperref}

\newtheorem{theorem}{Theorem}[section]
\newtheorem{lemma}[theorem]{Lemma}
\newtheorem{proposition}[theorem]{Proposition}
\newtheorem{definition}[theorem]{Definition}
\newtheorem{remark}[theorem]{Remark}

\newcommand{\N}{\mathbb N}
\newcommand{\WIN}{\mathrm{WIN}}
\newcommand{\LOSS}{\mathrm{LOSS}}
\newcommand{\DRAW}{\mathrm{DRAW}}

\title{A Binary Normal Form and an Audited Proof Attempt for the Two-Player $3n\mathbin{\pm}1$ Game}
\author{Grisha Pochuev\\
  \small\href{mailto:n_854@mail.ru}{n\_854@mail.ru}\\
  \small\url{https://github.com/Grisha-Pochuev/3n-plus-minus-1-game}\\
  \small\href{https://doi.org/10.5281/zenodo.21844684}%
  {doi:10.5281/zenodo.21844684}}
\date{Audited repair draft, 9 August 2026}

\begin{document}
\maketitle

\begin{center}
\small This preprint is distributed under the Creative Commons Attribution
4.0 International license (CC BY 4.0).
\end{center}

\begin{abstract}
We consider the impartial two-player game on positive odd integers in which a
move replaces $n>1$ by the odd part of either $3n-1$ or $3n+1$, and the player
who reaches $1$ wins. Arbitrary play need not terminate. We give an exact
binary normal form with one expanding and one strictly contracting branch and
develop a large collection of source, factor, and finite proof-height
normalizations. An external audit by Ingo Alth\"ofer identified a genuine gap
in the previous global entry argument: a decrease of the canonical odd
coefficient was incorrectly used as a decrease of its coefficient source,
and an arbitrary factorful exponent-one family was not covered by the reverse
factor lemma. The corrected supplement proves a universal two-level arithmetic
normalization for that missing family using Sections 79--81, but the returned
source can increase. The remaining provenance/rank attachment lemma is stated
explicitly in corrected Section 137. Consequently the global no-draw theorem
is \emph{not claimed as proved in this repair draft}. The computational,
symbolic-certificate, and Lean artifacts remain supporting checks with the
narrow trust boundary stated in the repository.
\end{abstract}

\section{Introduction}

The game studied here appears as Conway's \emph{Beans-Don't-Talk} game in
Guy's accounts \cite{Guy1986,Guy1996} and more recently as the $3n{+}{-}1$
game of Alth\"ofer, Hartisch, and Zipproth
\cite{AlthoferHartischZipproth2024}. Its similarity to Collatz iteration is
superficial in one important respect: the two signs are strategic choices by
alternating players.

The specific problem addressed here is the two-player problem stated on
Alth\"ofer's Collatz-problems page: prove that every odd starting value reaches
$1$ when both players act optimally \cite{AlthoferProblemPage}. The formulation
on that page and the game rules used below agree exactly.

For the player to move, a position is $\WIN$ if some move leads to a $\LOSS$
position, and is $\LOSS$ if every move leads to a $\WIN$ position. Because the
game graph is infinite, these two inductive classes need not exhaust all
positions. We call every remaining position $\DRAW$. Thus finite-game backward
induction cannot be used without a separate well-foundedness argument.

Indeed arbitrary play may cycle; for example $5\to7\to5$ is possible. The
main result concerns optimal play and is therefore strictly stronger than the
existence of some terminating choice and different from termination under all
choices.

\paragraph{Target theorem (open).}\label{thm:main}
The target statement is that for every positive odd starting integer, optimal
play reaches $1$ after finitely many moves, equivalently that the game has no
$\DRAW$ positions. The external audit described above shows that the current
manuscript does not yet prove this statement. Corrected Section 137 of the
supplement isolates the remaining attachment lemma exactly.

\paragraph{Independent verification package.}
The complete proof supplement and the independently runnable verification
artifacts are maintained at
\begin{center}
\url{https://github.com/Grisha-Pochuev/3n-plus-minus-1-game}.
\end{center}
An external reader can clone the repository and run \texttt{python audit.py}
to check the documented arithmetic regressions, finite proof certificates,
and the conditional global-routing assembly. The same repository contains the
negative tests and Lean-checked certificate metatheory, together with an
explicit account of what remains human-verified.

This manuscript presents the final dependency chain. Complete local case
proofs are contained in the numbered proof supplement
\cite{PochuevRepository}; section numbers below refer to that supplement.
The supplement, executable checks, certificate inventory, and formalization
sources are openly available at
\url{https://github.com/Grisha-Pochuev/3n-plus-minus-1-game}.

\section{Rules and binary normal form}

For a positive integer $x$, write $\operatorname{odd}(x)$ for the result of
dividing $x$ by the largest power of two that divides it. From an odd $n>1$
the two children are
\[
  \operatorname{odd}(3n-1),\qquad \operatorname{odd}(3n+1).
\]
The state $1$ is terminal and losing for the player to move.

Write $n=2m+1$ and define
\[
 F(m)=\left\lceil\frac{3m}{2}\right\rceil.
\]
Let $R(m)$ be the integer obtained by deleting the maximal alternating suffix
of the binary expansion of $m$. For instance,
$R(1001_2)=10_2$ and $R(10101_2)=0$.

\begin{proposition}[First normal form]\label{prop:first-normal}
For $m>0$, the two children in $m$-coordinates are exactly
\[
   F(m),\qquad F(R(m)).
\]
\end{proposition}

\begin{proof}
We have $3n-1=2(3m+1)$ and $3n+1=2(3m+2)$. Exactly one of
$3m+1,3m+2$ is odd. Splitting on the final bit of $m$ gives $F(m)$ from
that expression. In the other expression, each additional division by two
crosses one alternating bit of $m$; the first repeated adjacent bit stops the
division. This is precisely deletion of the maximal alternating suffix. The
four prefix/final-bit cases are given in the supplement's normal-form proof.
\end{proof}

Every child in Proposition~\ref{prop:first-normal} lies in the image of $F$.
Representing $F(q)$ by $q$ gives the conjugated game
\[
 A(q)=F(q)=\left\lceil\frac{3q}{2}\right\rceil,
 \qquad B(q)=R(F(q)),
\]
with terminal state $q=0$.

\begin{proposition}[Contracting branch]\label{prop:B-contracts}
For every $q>0$, $B(q)<q$.
\end{proposition}

\begin{proof}
Deleting a nonempty binary suffix gives
\[
B(q)\le \left\lfloor\frac{F(q)}2\right\rfloor
\le \left\lfloor\frac{3q+1}{4}\right\rfloor<q
\]
for $q\ge2$, while $B(1)=0$ directly.
\end{proof}

The difficulty is that $A$ expands and the length of the suffix deleted by
$R$ is unbounded. In particular no fixed residue modulus determines $R$.

\section{Finite proof tokens}

The ordinary inductive outcome construction gives every certified $\WIN$ or
$\LOSS$ position a finite canonical proof height. Terminal state $0$ has
height zero. A $\WIN$ certificate selects a lower-height $\LOSS$ child; a
$\LOSS$ certificate contains both lower-height $\WIN$ children.

During the draw analysis we retain the exact certified endpoints used by
these finite proofs. A marked configuration therefore carries a finite
multiset $M$ of proof heights. These are \emph{tokens}: their identities and
source provenance are retained, not merely their largest height or their
cardinality.

\begin{lemma}[Token multiset rank; Supplement Section 129]
For a finite multiset $M$ of proof heights define
\[
  \Phi(M)=\mathop{\#}_{h\in M}\omega^h,
\]
where $\#$ denotes natural ordinal sum. Replacing one token of height $h$ by
finitely many certified descendants of heights strictly below $h$ strictly
decreases $\Phi$.
\end{lemma}

The use of a multiset, rather than a scalar maximum, is essential: a routing
step may split one obligation into several lower ones. Equality of $\Phi$
alone is not used to identify configurations. The transition inventory proves
that every token-changing edge is strict, hence an equal-$\Phi$ route
preserves the actual token multiset.

\section{Marked factor routing}

Long constant binary tails admit coordinates
\[
 Q_D^g(C)=C2^D-g,
 \qquad g\in\{0,1\},
\]
with $C$ odd. The minimum-source analysis produces marked factor scans in
these coordinates. Their role is to keep a finite outcome witness attached
while the arithmetic route is normalized.

\begin{lemma}[Factor attachment; Supplement Sections 130--135]
Every marked factor scan retains its incoming finite proof tokens. For a
marked source $b=Q_D^g(C)$ with $D\ge3$, let $y$ be its marked $\LOSS$ child
and $q$ its marked $\WIN$ child. There is a common $\WIN$ child $r$ with
\[
  h(r)\le h(y)-1.
\]
Before an inner branch is followed, the scan replaces the carried occurrence
of $y$ by $r$. Every exit retains this lower token and its exact source
attachment. The remaining rows $D=1,2$ are exact lifts over an already marked
token and create no unranked fresh witness.
\end{lemma}

This statement closes a common logical gap: reaching a later finite endpoint
is insufficient unless it is comparable with the witness already used in the
outer rank. Sections 130--135 prove precisely that comparison.

\section{Well-founded fibres}

We use the following standard projection principle in a form adapted to the
marked relation.

\begin{lemma}[Well-founded fibre lemma; Supplement Section 132]
Let a transition relation carry a projection into a well-founded order.
Suppose the projection never increases, every projection-changing edge is
strictly decreasing, and the transition relation restricted to each fixed
projection fibre is well-founded. Then the full transition relation is
well-founded.
\end{lemma}

The proof is well-founded induction on the projection, followed by induction
inside the current fibre. The formulation permits inner finite heights that
are incomparable between different routes; only the retained outer projection
must be monotone.

\begin{lemma}[Typed fixed-fibre normalizer; corrected Supplement Section 136]
Once a provenance-preserving typed entry has supplied the retained source
anchor and certified proof tokens required by Sections 91--135, the
corresponding marked obligation/factor routing is well-founded in that fixed
retained fibre.
\end{lemma}

\begin{proof}[Scope of the repaired statement]
The retained arithmetic anchor is not the temporary obligation cursor.
Canonical $A$-continuations may increase the cursor and are treated as bounded
routing macros; they do not update the retained numerical rank. Strict source
changes and finite proof-token replacements are admitted only where the cited
lemmas prove the corresponding comparison. Sections 93--135 then control the
bounded $A$-streaks, valuation-two witnesses, factor forks, and high-return
cross-transitions for these typed entries. The missing issue is upstream: an
arbitrary factorful exponent-one DRAW has not yet been proved to acquire such
an entry without increasing the retained projection.
\end{proof}

\section{Audited entry gap and corrected arithmetic normalization}

The previous version asserted an exhaustive entry trichotomy. That assertion
was too strong. At a constant-tail boundary, Section 15 proves a decrease of
the canonical odd coefficient, not necessarily a decrease of the coefficient
\emph{source}. For example, $s=1$, $J(1)=5$, $a=15$, and $\epsilon=1$
give the common child $R(44)=22$, whose canonical coefficient is $11<15$ but
whose source is $3>1$ because $11=J(3)$.

The omitted arithmetic family is
\[
 Q_1^\epsilon(3^kJ(s)),\qquad k>0.
\]
The reverse-factor lemma applies only at exponent at least two. Corrected
Supplement Section 137 instead applies the universal two-level factor scan of
Sections 79--81. If the displayed state is DRAW, a DRAW occurs among the raw
and signed exits at the current or next factor level. Writing
\[
 3^{i+1}J(s)+1-2\epsilon=2^vJ(t),
\]
the signed exit is exactly $Q_v^{1-\epsilon}(J(t))$; for $v\ge2$ its raw
sibling is $Q_{v-1}^{1-\epsilon}(J(t))$, and for $v=1$ a positive raw DRAW
has coefficient source strictly below $t$.

This repairs the arithmetic exhaustiveness defect but does not prove that
$t<s$. In general the returned source may be larger. The remaining statement
is therefore a provenance statement, not another arithmetic factorization.

\paragraph{Remaining attachment lemma.}
Let $s$ be the globally least coefficient source of a DRAW and suppose
$Q_1^\epsilon(3^kJ(s))$ is DRAW with $k>0$. Starting from the two-level
normalization above, one must prove that an outcome-compatible finite
continuation either reaches an actual DRAW with retained source below $s$,
performs a certified strict replacement of an already carried finite proof
token, or enters the typed Section 136 normalizer while preserving the old
retained source/token projection. A temporary returned source $t$ may be
routing data but may not silently replace the retained anchor $s$.

For $s>0$ the ordinary state $s$ itself is finite, since its own coefficient
source is at most $(s-1)/6<s$. Sections 54, 57, and 61--64 therefore provide
finite proof-tree provenance that may be used in a completion. The source-zero
case is normalized separately in Sections 71--73.

\section{Current theorem status}

Because the attachment lemma above is still open, the no-$\DRAW$ theorem is
open in this manuscript. Sections 129 and 132 remain valid order-theoretic
lemmas, and corrected Section 136 remains valid for typed entries. If the
attachment lemma is proved, the previous final well-foundedness assembly can
be reinstated: strict retained-source and proof-token transitions cannot occur
infinitely, and an equal retained-rank path is trapped in the typed
well-founded normalizer. Until then, no such global conclusion is asserted.

\section{Verification status and reproducibility}

The repository distinguishes four kinds of evidence.

First, the mathematical supplement labels each claim as proved,
computationally verified, conjectural, disproved, or an unverified lead.
Second, standard-library Python code checks the exact normal form, descent
identities, routing arithmetic, and bounded retrograde soundness. These finite
runs are regression evidence and are not used as the proof of the infinite
theorem.

Third, the repository contains a finite declarative inventory of the global
routing assembly. It has 15 control states, 46 transition schemas, and 12
guard partitions, including symbolic unbounded valuation and tail-length
intervals. An independent standard-library Python checker verifies source
integrity, complete coverage of the \emph{declared} cases, one rule per case,
lexicographic reset discipline, and acyclicity of the 17 equal-rank control
edges. The remaining 29 schemas have an explicit strict rank component.

The checker prints the status \texttt{CONDITIONAL\_MACHINE\_CHECK}. Four
obligations deliberately remain on the human side of the trust boundary:
\begin{enumerate}
  \item the universal arithmetic and coordinate identities attached to each
        macro schema;
  \item refinement of the declared semantic guards to every legal game case;
  \item outcome compatibility of the selected DRAW continuation and the
        separately carried finite tokens;
  \item finite productivity of every macro-normalization.
\end{enumerate}
Thus acceptance of the JSON certificate validates only the declared typed
size-change assembly. Corrected Section 137 leaves one semantic attachment
obligation open, so certificate acceptance is not evidence for the global
no-$\DRAW$ theorem.

Fourth, a pinned Lean 4 project proves the general metatheory used by the
certificate: the lexicographic product of well-founded relations is
well-founded, a relation mapped to strict decreases in such a rank is
well-founded, and a well-founded certified relation admits no infinite route.
The Lean sources use no problem-specific axioms and contain no
\texttt{sorry} or \texttt{admit}. Lean does not yet prove the concrete
multiset rank or the four certificate-to-game refinement obligations, and the
coverage table states this limitation explicitly.

This separation is part of the result: it permits an auditor to reproduce the
computations without confusing them with the well-foundedness proof, and to
measure future formalization progress without importing the desired theorem
as an axiom.

From a clean checkout, the complete locally machine-checkable audit is:
\begin{verbatim}
python audit.py
\end{verbatim}
It runs the complete unit and regression test suite, checks the documented finite
arithmetic range, validates the global assembly certificate, generates a
finite outcome proof DAG, and validates that DAG with a separate checker. The
global certificate alone is checked by
\begin{verbatim}
python scripts/verify_global_certificate.py
\end{verbatim}
The Lean metatheory is checked independently by \texttt{lake build} in the
\texttt{formal} directory. Exact commands, expected output, negative tests,
and the human audit route are recorded in the repository's
\texttt{START\_HERE.md} and \texttt{AUDIT.md} files.

\section{Known proof hazards}

The argument does not propagate a least draw indefinitely along the expanding
branch, infer a theorem from a search depth, use a fixed modulus for an
unbounded suffix, or infer winning play from a merely terminating one-player
strategy. It also does not replace retained witnesses by incomparable later
witnesses. Most importantly, the fixed-fibre proof checks transitions between
the obligation and factor subsystems; their separate termination would not by
itself imply termination of their union.

\section{Conclusion}

The binary normal form and the local factor/source identities survive the
external audit. The audit nevertheless invalidates the previous global proof:
the arbitrary factorful exponent-one family was not attached to the retained
well-founded rank. The corrected supplement now gives a complete two-level
arithmetic normalization of that family and isolates the remaining
provenance/rank attachment lemma. The global termination statement remains
open until that lemma is proved and independently re-audited.

\section*{Data, code, and correspondence}

The proof supplement, source code, tests, certificate files, Lean sources, and
reproducibility records are available at
\begin{center}
\url{https://github.com/Grisha-Pochuev/3n-plus-minus-1-game}.
\end{center}
The repository is the authoritative location for the evolving verification
artifacts. The released article and source are permanently archived at
\begin{center}
\href{https://doi.org/10.5281/zenodo.21844684}%
{doi:10.5281/zenodo.21844684}.
\end{center}
Questions, counterexamples, and audit reports may be sent to
\href{mailto:n_854@mail.ru}{n\_854@mail.ru}.

\bibliographystyle{plain}
\bibliography{references}

\end{document}
